% Pandoc передаёт выбранный язык через параметры класса документа.
\usepackage[provide=*,shorthands=off]{babel}
% Шрифты выбираются по именам семейств, независимо от расположения проекта.
\IfFontExistsTF{Segoe UI Emoji}{%
  \directlua{
    local extra = {"Segoe UI Symbol:mode=harf;", "Segoe UI Emoji:mode=harf;"}
    local windows = os.getenv("WINDIR")
    if windows then
      table.insert(extra, "[" .. windows:gsub(string.char(92), "/") .. "/Fonts/tahoma.ttf]:mode=harf;")
    end
    luaotfload.add_fallback("bookglyphs", extra)
  }%
  \defaultfontfeatures+[\rmfamily,\sffamily,\ttfamily]{RawFeature={fallback=bookglyphs}}%
}{}
\IfFontExistsTF{DejaVu Serif}
  {\setmainfont{DejaVu Serif}}
  {\setmainfont{Times New Roman}}
\IfFontExistsTF{DejaVu Sans}
  {\setsansfont{DejaVu Sans}}
  {\setsansfont{Arial}}
\IfFontExistsTF{DejaVu Sans Mono}
  {\setmonofont[Ligatures=NoCommon]{DejaVu Sans Mono}}
  {\setmonofont[Ligatures=NoCommon]{Consolas}}

\usepackage{fvextra}
\usepackage{needspace}
\usepackage{bookmark}
\hypersetup{pdflang={\BCPdata{main.tag}}}
\RecustomVerbatimEnvironment{verbatim}{Verbatim}{breaklines=true,breakanywhere=true}
\ifcsname Highlighting\endcsname
  \RecustomVerbatimEnvironment{Highlighting}{Verbatim}
    {commandchars=\\\{\},breaklines=true,breakanywhere=true}
  % Учитываем возможный перенос и внутри длинной строки одного цвета.
  \newcommand{\BookWrapToken}[1]{%
    \ifcsdef{#1Tok}{%
      \csletcs{BookOriginal#1Tok}{#1Tok}%
      \csdef{#1Tok}##1{\csuse{BookOriginal#1Tok}{\FancyVerbBreakStart##1\FancyVerbBreakStop}}%
    }{}%
  }
  \AtBeginEnvironment{Highlighting}{\forcsvlist{\BookWrapToken}{Alert,Annotation,Attribute,BaseN,BuiltIn,Char,Comment,CommentVar,Constant,ControlFlow,DataType,DecVal,Documentation,Error,Extension,Float,Function,Import,Information,Keyword,Normal,Operator,Other,Preprocessor,RegionMarker,SpecialChar,SpecialString,String,Variable,VerbatimString,Warning}}
\fi
\pdfstringdefDisableCommands{\def\allowbreak{}}

\makeatletter
\renewcommand{\@pnumwidth}{3em}
\renewcommand{\@tocrmarg}{4em}
\renewcommand{\l@subsection}{\@dottedtocline{2}{1.5em}{3.6em}}
\renewcommand{\l@subsubsection}{\@dottedtocline{3}{4.5em}{5.3em}}
\newcommand{\BookDeepHeading}[3]{%
  \phantomsection
  \ifnum#2=1
    \edef\BookBookmark{\noexpand\bookmark[level=#1,dest={\@currentHref}]}%
    \BookBookmark{#3}%
  \fi
  \@startsection{paragraph}{4}{\z@}{1.5ex}{.6ex}{\normalfont\normalsize\bfseries}*{#3}%
}
\newbox\BookHeadingBox
\newcommand{\BookTableHeadings}[1]{%
  \par\addvspace{.65\baselineskip}%
  \setbox\BookHeadingBox=\vbox{\hsize=\linewidth\@nobreakfalse #1\par}%
}
\newcommand{\BookHeadings}[1]{%
  \BookTableHeadings{#1}%
  \Needspace{\dimexpr\ht\BookHeadingBox+\dp\BookHeadingBox+2\baselineskip\relax}%
  \unvbox\BookHeadingBox
  \@afterindentfalse\@afterheading
}
% Первая страница длинной таблицы должна содержать и ведущие в неё заголовки.
\AddToHook{package/longtable/after}{%
  \AddToHook{env/longtable/begin}{\ifx\LTcaptype\@empty\stepcounter{LT@tables}\fi}%
  \AddToHook{cmd/LT@start/before}{%
    \ifvoid\BookHeadingBox\else
      \setbox\LT@firsthead=\vbox{%
        \unvbox\BookHeadingBox
        \vskip.65em\relax
        \ifvoid\LT@firsthead\copy\LT@head\else\unvbox\LT@firsthead\fi}%
      \global\setbox\BookHeadingBox=\box\voidb@x
    \fi
  }%
}
\makeatother
